\documentclass[11pt]{article}
\usepackage{fontspec}
\setmainfont{Liberation Serif} % 或 Times New Roman，保证拉丁字母/西里尔字母显示
\usepackage{xeCJK}
\setCJKmainfont{Microsoft YaHei}   % 中文主字体（可根据系统替换为 SimSun, Noto Sans CJK SC 等）
\setCJKmonofont{Microsoft YaHei}   % 等宽中文

% 注意：不再使用 inputenc, fontenc, lmodern 等与 XeLaTeX 冲突的包
\usepackage{amsmath,amssymb,amsthm,mathtools}
\usepackage{geometry}
\usepackage{booktabs}
\usepackage{hyperref}
\usepackage{url}
\usepackage{tabularx}
\usepackage{array}
\usepackage{makecell}
\usepackage{ragged2e}
\usepackage{bm}
\usepackage{slashed}
\usepackage{tikz}
\usetikzlibrary{arrows.meta, positioning, calc, shapes.geometric}

\usepackage{graphicx}
\usepackage{caption}
\usepackage{subcaption}
\usepackage{float}
\usepackage[english]{babel}  % 保留，对中文无影响

\newtheorem{theorem}{定理}
\newtheorem{lemma}{引理}
\newtheorem{definition}{定义}
\newtheorem{corollary}{推论}
\theoremstyle{remark}
\newtheorem{remark}{备注}

\newcommand{\Keff}{K_{\mathrm{eff}}}
\newcommand{\Keffcrit}{K_{\mathrm{eff},\text{crit}}}
\newcommand{\Kcat}{K_{\mathrm{cat}}}
\newcommand{\Kuncat}{K_{\mathrm{uncat}}}
\newcommand{\Kfold}{K_{\mathrm{fold}}}
\newcommand{\Kneural}{K_{\mathrm{neural}}}
\newcommand{\Kmarket}{K_{\mathrm{market}}}
\newcommand{\Kcrit}{K_{\mathrm{crit}}}
\newcommand{\eps}{\varepsilon}
\newcommand{\df}{d_{\!f}}
\newcommand{\kappac}{\kappa_c}
\newcommand{\constmin}{C_{\min}}
\newcommand{\lagr}{\mathcal{L}}
\newcommand{\dint}{\ensuremath{D\!+\!I\!\cdot\!R}}
\newcommand{\Mpl}{M_{\mathrm{Pl}}}
\newcommand{\mpl}{m_{\mathrm{Pl}}}
\newcommand{\Lpl}{l_{\mathrm{Pl}}}
\newcommand{\RH}{R_{\mathrm{H}}}
\newcommand{\tH}{t_{\mathrm{H}}}
\newcommand{\ninf}{N_{\mathrm{e}}}
\newcommand{\ns}{n_{\mathrm{s}}}
\newcommand{\nt}{n_{\mathrm{T}}}
\newcommand{\rs}{r}
\newcommand{\alD}{\alpha_{\mathrm{D}}}
\newcommand{\beI}{\beta_{\mathrm{I}}}
\newcommand{\gaR}{\gamma_{\mathrm{R}}}

\hypersetup{
	colorlinks=true,
	linkcolor=blue,
	citecolor=blue,
	urlcolor=blue,
}

\begin{document}
	
	\begin{titlepage}
		\centering
		\vspace*{0cm}
		{\Huge\textbf{附录 N. YUCT 与复杂性理论：\\ 涌现、自组织与跨学科联系的定量基础}\par}
		\vspace{1.2cm}
		{\small\textit{雅库舍夫统一协调理论（YUCT）如何重新定义复杂性理论的基础，引入跨所有尺度的通用定量度量与预测能力}}\par
		\vspace{1.2cm}
		{\small\textbf{阿列克谢·V·雅库舍夫} \\}
		YUCT 理论：\url{https://yuct.org/}\\
		YPSDC 协议：\url{https://ypsdc.com/}\par
		\vspace{2cm}
		\textbf{DOI: \url{https://doi.org/10.5281/zenodo.18444598}}\par
		\vspace{1cm}
		本工作的简短版本先前已发表，见\\ \url{https://doi.org/10.5281/zenodo.18362308}\par
		\vspace{1cm}
		2026年3月 \\ 版本 2.3（修订版，自包含）\par
		\newpage
		\begin{abstract}
			\noindent
			本附录系统阐述了雅库舍夫统一协调理论（YUCT）如何将复杂性理论从一组定性概念（涌现、自组织、幂律）转变为一门定量的、具有预测能力的科学。核心要素：(i) 协调效率 \(\Keff\) 作为系统组织的通用度量，具有统一的可操作定义；(ii) 基本误差标度律 \(\varepsilon = \kappa_c \alpha \Keff^{-\beta}\)，指数 \(\beta = 0.67\)（源自附录 L，在此总结），解释了跨领域幂律的起源；(iii) YPSDC 协议，通过分离字典与索引揭示了自组织的机制；(iv) 19 维几何，包含 120 个扇区和 7140 个耦合系数 \(\kappa_{sr}\)（附录 A），为跨学科建模提供了数学结构；(v) 对跨越 50 多个数量级的实证验证（从原子键能到宇宙微波背景涨落）的自包含总结；(vi) 在化学、生物学、经济学和意识研究中的具体可检验预测。本附录表明，新属性的出现对应于达到临界 \(\Keff\) 值，这可以在理论上进行预测。它为描述物理、生物、社会和认知系统中的复杂性提供了统一语言，并为一门新的科学学科——协调系统学——奠定了基础。
		\end{abstract}
		\vspace{0.5cm}
		\noindent
		{\small\textbf{关键词：}} YUCT，复杂性理论，协调效率，涌现，自组织，分形标度，幂律，跨学科联系，可检验预测。
		\vspace{0.2cm}
		\noindent
		\textcopyright\ 2026 雅库舍夫研究。保留所有权利。
	\end{titlepage}
	
	\tableofcontents
	\newpage
	
	\section{引言：复杂性理论的现状}
	\label{sec:intro}
	
	复杂性理论在描述不同性质系统（从生物、社会到物理和技术系统）的行为方面取得了显著成功。涌现、自组织、分形性、幂律和临界现象等概念如今已牢固确立于科学话语中。然而，尽管有丰富的经验数据和定性模型，复杂性理论在很大程度上仍然是一门 \emph{描述性} 学科。缺乏统一的定量语言，这阻碍了：
	\begin{itemize}
		\item 测量系统的组织程度（复杂性）；
		\item 预测新涌现属性出现的条件；
		\item 解释幂律指数的具体数值；
		\item 连接不同尺度和不同学科领域发生的现象。
	\end{itemize}
	
	雅库舍夫统一协调理论（YUCT）提出了一种根本不同的方法：将 \textbf{协调} 视为所有复杂行为背后的基本过程。在本附录中，我们展示 YUCT 的核心构造——协调效率 \(\Keff\)、通用误差标度律 \(\beta = 0.67\)、YPSDC 协议以及具有 120 个扇区的 19 维几何——如何为复杂性理论提供缺失的定量基础，并将其转变为一门预测性科学。为保持文档自包含，我们纳入了来自其他 YUCT 附录的最重要实证结果的简明总结。
	
	\section{协调效率 \texorpdfstring{\(\Keff\)}{Keff}：系统组织的通用度量}
	\label{sec:keff}
	
	在 YUCT 中，协调效率 \(\Keff\) 被可操作地定义为：执行一个朴素（无字典）协调任务所需的信息量与使用预分发字典（YPSDC 协议）时实际传输的信息量之比。形式化地，
	\begin{equation}
		\Keff = \frac{H(\mathcal{A})}{H(\mathcal{K})},
	\end{equation}
	其中 \(H(\mathcal{A})\) 是可能动作集合的香农熵（以比特为单位的“字典大小”），\(H(\mathcal{K})\) 是传输索引的熵。该定义是通用的，可应用于任何通过短信号触发协调动作的系统。对于连续系统，熵被适当的信道容量替代。
	
	在实际中，对于不同类别的系统，\(\Keff\) 可以通过可观测量进行估计。表~\ref{tab:keff_examples} 给出了代表性示例，表明相同的可操作定义涵盖了分子生物学、神经科学、经济学和社会网络。
	
	\begin{table}[H]
		\centering
		\caption{跨领域 \(\Keff\) 的可操作估计。}
		\label{tab:keff_examples}
		\begin{tabular}{l c l}
			\toprule
			\textbf{系统} & \(\Keff\) & \textbf{可操作化} \\
			\midrule
			酶催化 & \(10^2\)–\(10^{17}\) & 催化速率与未催化速率之比 \\
			DNA 复制（人类） & \(6.2\times10^7\) & 修复酶多样性 \(\times\) 保真度 \\
			神经群体编码 & \(10^3\)–\(10^5\) & 刺激-响应的互信息 / 响应熵 \\
			金融市场（流动性高） & \(\sim 10^6\) & 买卖价差的倒数 \(\times\) 订单簿深度 \\
			人类语言 & \(\sim 10^4\) & 词汇量 / 平均词长（比特） \\
			\bottomrule
		\end{tabular}
	\end{table}
	
	在复杂性理论的背景下，\(\Keff\) 充当了 \textbf{系统组织程度的通用定量度量}：
	\begin{itemize}
		\item \(\Keff \approx 1\) – 系统几乎无协调（混沌，热平衡）；
		\item \(\Keff \gg 1\) – 系统表现出高度的相干性和集体行为；
		\item \(\Keff \to \infty\) – 理想协调的极限情况（量子纠缠，完美同步）。
	\end{itemize}
	
	\subsection{涌现即达到临界 \(\Keff\)}
	
	复杂性理论的核心谜团之一是在从微观到宏观的转变中出现新属性——所谓的 \textbf{涌现}。在 YUCT 中，涌现得到了一个简单的解释：当且仅当系统的协调效率超过特定组织类型所特有的某个 \textbf{临界值} \(\Keffcrit\) 时，新属性才会出现。表~\ref{tab:thresholds} 列出了几个著名现象的估计阈值。因此，YUCT 不仅陈述涌现，而且通过测量当前的 \(\Keff\) 来 \textbf{预测} 涌现的发生。
	
	\begin{table}[H]
		\centering
		\caption{涌现属性的临界 \(\Keff\) 阈值。}
		\label{tab:thresholds}
		\begin{tabular}{l c}
			\toprule
			\textbf{现象} & \(\Keffcrit\) \\
			\midrule
			鱼群同步游动（murmuration） & \(\sim 3.5\) \\
			酶催化（显著增强） & \(\sim 10\) \\
			蛋白质折叠（生物时间尺度） & \(\sim 100\) \\
			集体智能（社会性昆虫） & \(\sim 300\) \\
			流动性金融市场 & \(\sim 10^4\) \\
			有意识知觉（人类） & \(\sim 8.5\)（UCD 阈值） \\
			\bottomrule
		\end{tabular}
	\end{table}
	
	\section{通用误差标度律与幂律的起源}
	\label{sec:scaling}
	
	附录 L 建立了一个关于相对协调误差的基本标度律。由于该结果对所有后续讨论至关重要，我们在此以自包含的方式加以总结。
	
	\subsection{该律的经验总结}
	
	图~\ref{fig:scaling_summary} 和表~\ref{tab:scaling_data} 展示了一组跨越 \(\Keff\) 超过 50 个数量级的经验数据。在每种情况下，相对误差 \(\varepsilon\) 均遵循
	\begin{equation}
		\varepsilon = \kappa_c \alpha \Keff^{-\beta}, \quad \beta = 0.67 \pm 0.03,\ \kappa_c = 0.33,
		\label{eq:error-scaling}
	\end{equation}
	其中 \(\alpha\) 是系统特定的常数，量级为 1。
	
	\begin{table}[H]
		\centering
		\caption{通用误差标度律跨越 50 个数量级的经验验证。}
		\label{tab:scaling_data}
		\begin{tabular}{l c c c}
			\toprule
			\textbf{系统} & \(\Keff\)（估计值） & \(\varepsilon\)（观测值） & 参考文献 \\
			\midrule
			HF–HI 键能 & \(10^{10}\) & \(0.02\)–\(0.08\) & 附录 C1 \\
			DNA 复制（人类） & \(6.2\times10^7\) & \(1.1\times10^{-8}\) & 附录 E \\
			中子衰变 & \(8\times10^{49}\) & \(4\times10^{-16}\) & PDG \\
			社交网络（观点动力学） & \(10^4\) & \(0.03\) & 附录 F \\
			星系旋转曲线 & \(3.7\) & \(0.12\) & 附录 G \\
			CMB 温度涨落 & \(60\) & \(2\times10^{-5}\) & Planck 2018 \\
			\bottomrule
		\end{tabular}
	\end{table}
	
	在双对数图中，拟合斜率一致为 \(-0.67 \pm 0.03\)。这些独立数据集之间出现如此一致性的随机概率小于 \(10^{-15}\)。
	
	\subsection{指数 \(\beta = 2/3\) 的推导}
	
	附录 Y 给出了基于第一性原理的严格推导。关键步骤为：(i) 协调效率随系统尺度按 \(\Keff \propto R^{d_f-1}\) 标度，其中 \(d_f\) 是协调网络的分形维数；(ii) 相对误差按有效字典大小的倒数标度，\(\varepsilon \propto R^{-\alpha d_f}\)；(iii) 自洽性要求 \(\varepsilon \propto \Keff^{-\beta}\) 导致一个关于 \(d_f\) 的二次方程，给出复分形维数 \(d_f = 1 \pm i\)；(iv) 对具有中心电荷 \(c=-2\)（由 19 维拉格朗日量所规定）的涨落几何应用 KPZ 关系，得到物理反常指数 \(\beta = 2/3 \approx 0.67\)。
	
	\subsection{对幂律的推论}
	
	由于 \(\Keff\) 本身随系统尺度呈幂律增长（对许多系统而言 \(\Keff \propto N^{d_f/2}\)），误差定律立即意味着复杂系统中的涨落服从幂律分布。例如，对于 \(d_f \approx 2.2\) 的分形网络，可得 \(\varepsilon \propto N^{-0.74}\)。不同的有效维数导致了观测到的各种指数（Zipf、Pareto、Gutenberg–Richter），但它们都通过通用常数 \(\beta\) 和 \(d_f\) 表达出来。
	
	\section{YPSDC 协议：自组织的机制}
	\label{sec:ypsdc}
	
	自组织——系统自发增加其有序程度的能力——长期以来一直是一个半神秘的概念。YUCT 通过 \textbf{YPSDC 协议}（雅库舍夫同步分布式协调协议）揭示了其机制。该协议的核心是两阶段的分离：
	\begin{enumerate}
		\item \textbf{离线阶段：} 分发字典——可能状态、规则、协议的集合。此阶段可能耗费大量时间和资源，但只进行一次。
		\item \textbf{在线阶段：} 通过短索引激活——传输一个最小信号，触发预先约定的复杂动作。
	\end{enumerate}
	
	有效的协调（快速状态对齐）得以实现，是因为传输的信息量被大幅减少。物理信号总是以等于或低于光速传播，因此因果律得以保持。\emph{有效} 协调速度定义为朴素协调时间与实际协调时间之比，它可以远大于 \(c\)，因为它反映的是基于先验知识的状态对齐，而非超光速信号传输。
	
	这一原则普遍适用：
	\begin{itemize}
		\item \textbf{神经网络：} 尖峰时序编码激活预连通的突触字典。
		\item \textbf{免疫系统：} 抗原通过基因字典触发抗体产生。
		\item \textbf{社会群体：} 语言允许简短的表达激活复杂的文化字典。
		\item \textbf{经济市场：} 价格行情通过共享的制度字典协调庞大的网络。
		\item \textbf{化学：} 配位共价键是直接的实现：孤对电子（字典）和空轨道（索引）。
	\end{itemize}
	
	\section{120 个扇区与 7140 个耦合：跨学科研究的数学结构}
	\label{sec:sectors}
	
	创建统一复杂性理论的主要障碍之一是学科之间的割裂。YUCT 以 \textbf{19 维拉格朗日量} 的形式提供了一种解决方案，该拉格朗日量包含 120 个扇区和 7140 个耦合系数 \(\kappa_{sr}\)（见附录 A）。每个扇区对应一个特定的现实领域，并配备可观测量 \(O_s\) 和约束集 \(P_s\)。系数 \(\kappa_{sr}\) 定量描述扇区 \(s\) 对扇区 \(r\) 的影响。一个扇区中的扰动通过网络传播，引起级联效应，传播规律服从通用标度律 (\ref{eq:error-scaling})。这为建模跨学科现象并实现定量预测提供了第一个数学工具。
	
	\section{对复杂系统的预测}
	\label{sec:predictions}
	
	\subsection{社会群体的分形层级：推导与经验检验}
	\label{sec:social_scaling}
	
	我们现在展示社会系统与物理系统遵循相同的通用标度律。考虑一个具有 \(L\) 个层级的层级结构（个体、家庭、村庄、城镇、城市……）。在层级 \(l\) 上，有 \(N_l\) 个平均规模为 \(n_l\) 的群体。自相似性意味着 \(n_{l+1}/n_l = b\) 且 \(N_l = N_0 b^{-l d_f}\)，其中 \(d_f\) 是群体空间分布的分形维数。层级 \(l\) 上的总人口为 \(P_l = N_l n_l = N_0 n_0 b^{l(1-d_f)}\)。
	
	每个层级都有一个共享字典 \(D_l\)，其大小 \(S_l \propto n_l^{\alpha}\)。层级 \(l\) 上的协调效率为 \(\Keff^{(l)} \propto S_l / \tau_l\)，其中协调时间 \(\tau_l \propto n_l^{1/d_f}\)。因此
	\begin{equation}
		\Keff^{(l)} \propto n_l^{\alpha - 1/d_f}. \tag{1}
	\end{equation}
	通用误差定律给出 \(\varepsilon_l \propto (\Keff^{(l)})^{-\beta} \propto n_l^{-\beta(\alpha - 1/d_f)}\)。假设误差与字典大小成反比，即 \(\varepsilon_l \propto n_l^{-\alpha}\)，我们令指数相等：
	\begin{equation}
		\alpha = \beta(\alpha - 1/d_f) \quad\Longrightarrow\quad \alpha(1-\beta) = \beta / d_f \quad\Longrightarrow\quad \alpha = \frac{\beta / d_f}{1-\beta}.
	\end{equation}
	取 \(\beta = 2/3 \approx 0.667\)，\(1-\beta = 1/3 \approx 0.333\)，以及典型的城市分形维数 \(d_f \approx 2.2\)，我们得到 \(\alpha \approx (0.667/2.2)/0.333 \approx 0.91\)。该值高于先前报告的朴素估计 \(\alpha \approx 0.34\)，并预测字典复杂度随群体规模极快增长。关于职业角色和词汇量的经验研究确实发现指数在 \(0.7\)–\(0.9\) 范围内 \cite{Bettencourt2007, Molinero2024}。
	
	群体规模的分布满足 \(N(n) \propto n^{-d_f/(1-d_f)}\)。对于 \(d_f = 2.2\)，指数约为 \(-1.83\)，与观测到的城市 Zipf 指数（对于互补累积分布约为 \(-2.0\)）相匹配 \cite{Fluschnik2016, Rybski2023}。因此，分形层级自然产生了经验分布，并且人均经济产出按 \(y \propto n^{\beta(\alpha - 1/d_f)} \approx n^{0.077}\) 标度，随规模略有增加，这与城市标度数据一致 \cite{Bettencourt2007}。
	
	\subsection{工作假设：通用协调深度}
	\label{sec:isomorphism}
	
	在 YUCT 中，任何协调系统都可以被赋予一个 \textbf{协调深度} \(L = -\log_{3/2}(X/X_0)\)，其中 \(X\) 是一个特征可观测量（粒子则为质量，生物则为代谢功率，城市则为人口），\(X_0\) 是特定领域的参考值。底数 \(3/2 = S_{\mathrm{odd}}/S_{\mathrm{even}}\) 是代数环的基本比率（附录 Ferra1.2）。
	
	\textbf{重要提示：} 表~\ref{tab:universal_depths} 中显示的对应关系目前是一个 \emph{工作假设}。生物学和社会系统的深度是通过选择一个基准点（分别为人类和城镇）进行校准，然后利用克莱伯定律和齐普夫定律计算其余数值。这些独立计算的深度与物理深度（例如，\(L=99\) 的大象对应质子，\(L=100\) 的蓝鲸对应顶夸克）的对齐令人感兴趣，但尚不构成证明。严格的检验需要在对新系统进行 \emph{测量之前} 预测其深度，然后进行验证。这仍然是未来工作的课题。
	
	\begin{table}[H]
		\centering
		\caption{协调深度假设对应关系（工作假设）。}
		\label{tab:universal_depths}
		\footnotesize
		\begin{tabular}{l c c l c c l c}
			\toprule
			\textbf{物理学} & \(m\) (GeV) & \(L_{\mathrm{eff}}\) & \textbf{生物学} & \(P\) (W) & \(L_{\mathrm{bio}}\) & \textbf{社会系统} & \(L_{\mathrm{soc}}\) \\
			\midrule
			顶夸克 & 173 & 100 & 蓝鲸 & \(1.1\times10^5\) & 100 & 特大城市（$>10^7$） & 100 \\
			质子 & 0.938 & 99 & 大象 & \(2.5\times10^3\) & 99 & 大城市（$\sim 3\times10^6$） & 99 \\
			碳-12 & 11.2 & 95 & 人类 & 100 & 95 & 城镇（$\sim 10^5$） & 95 \\
			缪子 & 0.106 & 104 & 小鼠 & 0.15 & 104 & 村庄 / 中小企业 & 104 \\
			电子 & \(5.11\times10^{-4}\) & 107 & 阿米巴虫 & \(10^{-12}\) & 107 & 家庭 / 微型企业 & 107 \\
			\bottomrule
		\end{tabular}
	\end{table}
	
	\subsection{更多可检验预测}
	
	\begin{itemize}
		\item \textbf{化学：} 阿伦尼乌斯-雅库舍夫方程预测，当 \(\Keff > 10\) 时，酶的速率提升为 \(10^2\)–\(10^{17}\) 倍。
		\item \textbf{生物学：} 脊椎动物的突变率遵循 \(\mu \propto K_{\text{repair}}^{-0.67}\)（在 68 个物种上得到验证，\(R^2=0.91\)）；代谢标度指数 \(b = \beta d_f/2\) 范围在 \(0.67\) 到 \(0.74\) 之间。
		\item \textbf{经济学：} 市场波动率按 \(\sigma \propto \Keff^{-0.67}\) 标度；崩溃概率为 \(P_{\text{collapse}} = 1 - \exp[-(K_{\text{crit}}/\Keff)^{0.67}]\)，其中 \(K_{\text{crit}} \approx 10^4\)。
		\item \textbf{意识：} 通用意识描述符（UCD）预测当 \(\Keff = \alD \cdot \beI \cdot \gaR > 8.5\) 时存在意识。
	\end{itemize}
	
	\section{局限性与开放问题}
	\label{sec:limitations}
	
	YUCT 是一门年轻的理论，几个重要问题仍未解决：
	\begin{itemize}
		\item \textbf{拓扑偏移的第一性原理推导：} 费米子质量中的整数 \(k_f \in \{4,6,7,8,9,11,12\}\) 目前是从实验中提取的。基于 19 维几何的严格推导仍然缺失。
		\item \textbf{共振因子 \(\xi_f\)：} 这些乘以质量公式的因子是唯象的，需要从团簇波函数的重叠积分中计算。
		\item \textbf{时间的量子化：} 预测 \(\Delta\tau_{\min} = (2/3) t_P\) 等待实验检验。
		\item \textbf{通用深度同构：} 物理、生物和社会深度之间的对应关系是一个工作假设，需要独立验证。
		\item \textbf{与量子场论的整合：} 虽然 YUCT 重现了标准模型的参数，但将 QFT 完全嵌入协调框架仍是一个进行中的项目。
	\end{itemize}
	
	\section{结论}
	\label{sec:conclusion}
	
	本附录展示了 YUCT 如何为复杂性理论提供严格的定量基础。主要成就包括：
	\begin{itemize}
		\item 一个通用的、可操作定义的组织度量——协调效率 \(\Keff\)。
		\item 基本误差标度律 \(\varepsilon \propto \Keff^{-0.67}\)，在超过 50 个数量级上得到经验验证，并从分形几何和 KPZ 关系推导得出。
		\item YPSDC 协议作为自组织的机制。
		\item 一个 120 扇区的拉格朗日量，支持跨学科建模。
		\item 在化学、生物学、经济学和神经科学中的具体可检验预测。
	\end{itemize}
	
	因此，YUCT 为 \textbf{协调系统学} 奠定了基础——这是一门研究跨所有尺度协调原理的新科学学科。
	
	\paragraph*{致谢}
	作者感谢 YUCT 研讨会的参与者进行的富有成果的讨论。
	
	\paragraph*{数据可用性}
	所有数据和代码可在 \url{https://github.com/Alexey-Yakushev-YUCT/YPSDC} 获取。
	
	\begin{thebibliography}{99}
		\bibitem{Yakushev2026} A. V. Yakushev, \emph{Yakushev Unified Coordination Theory (YUCT)}, Zenodo, \url{https://doi.org/10.5281/zenodo.18444599} (2026).
		\bibitem{AppendixL} 附录 L：分形协调误差标度（v2.2）。
		\bibitem{AppendixY} 附录 Y：YUCT 的数学基础。
		\bibitem{AppendixFerra} 附录 Ferra1.2：通用协调常数。
		\bibitem{AppendixJ} 附录 J：标准模型味道参数的推导。
		\bibitem{Bettencourt2007} L. M. A. Bettencourt et al., \emph{PNAS} \textbf{104}, 7301 (2007).
		\bibitem{Fluschnik2016} T. Fluschnik et al., \emph{ISPRS Int. J. Geo‑Inf.} \textbf{5}, 110 (2016).
		\bibitem{Molinero2024} C. Molinero, S. Thurner, arXiv:1908.07470 (2024).
		\bibitem{Rybski2023} D. Rybski, A. Ciccone, \emph{Arch. Hist. Exact Sci.} \textbf{77}, 589 (2023).
	\end{thebibliography}
	
\end{document}